\section*{Introduction}
Linear algebraic groups lie exactly at the intersection of algebraic geometry and representation theory. Unlike Lie groups, which are modeled on smooth manifolds, algebraic groups are rooted in the framework of algebraic varieties. The primary aim of this thesis is to explore this geometric-algebraic interplay, specifically focusing on connected reductive algebraic groups over an algebraically closed field of characteristic zero, and culminating in the geometric realization of their representations via the Borel-Weil theorem. By studying how the underlying variety structures dictate group actions and representations, we describe a link between the topology of flag varieties and the algebraic classification of weights.

To build this framework, this project first establishes the required language of algebraic geometry, moving from affine and projective varieties to the geometry of Grassmannians and flag varieties. By utilizing the Plücker embedding, we formalize the projective structure of these spaces, establishing a concrete frame for subsequent geometric constructions. This geometric foundation immediately interfaces with the theory of algebraic groups, where group operations are regular maps. We analyze the structural anatomy of these groups, focusing on maximal tori, Borel subgroups, and the adjoint representation on the associated Lie algebra, to isolate the class of reductive groups. A key geometric point of this structural theory is the realization of the quotient space $G/B$ as a smooth projective variety, which naturally generalizes the classical flag variety and serves as the central geometric object of the entire theory.

The narrative then transitions into the classical representation theory of these groups, examining how the action of a maximal torus decomposes a rational representation into weight spaces. By introducing the combinatorial machinery of root systems and Weyl groups, we arrive at the Highest Weight Theorem. This foundational result provides a complete classification of irreducible finite-dimensional representations, demonstrating that each is uniquely characterized by a dominant integral weight. While this algebraic classification is highly effective, it leaves open the question of whether these representations can be realized intrinsically through the geometry of the group itself.

This question finds its resolution in the Borel-Weil theorem, which serves as the final goal of this project. By assigning an algebraic line bundle over the flag variety $G/B$ to each character of a maximal torus, we allow the geometric study of representations. The theorem demonstrates that the space of global sections of these line bundles is precisely the duals of the irreducible highest-weight representations.

For the rest of this project, let $K$ denote an algebraically closed field, which we think of as $\CC$.
